% ============================================================
%  Informe Laboratorio 3 - Sensores y librerías
%  Compilar con: pdflatex informe_lab3.tex  (dos veces)
% ============================================================
\documentclass[11pt,a4paper]{article}

\usepackage[utf8]{inputenc}
\usepackage[T1]{fontenc}
\usepackage[spanish,es-noquoting]{babel}
\usepackage{lmodern}
\usepackage[margin=2.2cm,top=1.9cm,bottom=1.9cm]{geometry}
\usepackage{xcolor}
\usepackage{enumitem}
\usepackage{listings}
\usepackage{titlesec}
\usepackage[colorlinks=true,urlcolor=blue!55!black]{hyperref}

\setlist[itemize]{topsep=2pt,itemsep=1.5pt,parsep=0pt,leftmargin=1.2em}
\setlength{\parskip}{3pt}
\pagestyle{empty}

\titleformat{\section}{\normalsize\bfseries}{}{0pt}{}
\titlespacing*{\section}{0pt}{9pt}{2pt}

% --- estilo de los listados de código ---
\lstset{
  language=C,
  basicstyle=\ttfamily\scriptsize,
  commentstyle=\color{black!55}\itshape,
  keywordstyle=\bfseries,
  showstringspaces=false,
  tabsize=2,
  frame=leftline,
  framerule=0.6pt,
  rulecolor=\color{black!40},
  xleftmargin=8pt,
  aboveskip=4pt,
  belowskip=2pt,
  literate={á}{{\'a}}1 {é}{{\'e}}1 {í}{{\'i}}1 {ó}{{\'o}}1 {ú}{{\'u}}1
           {ñ}{{\~n}}1 {¿}{{?`}}1 {→}{{$\to$}}1
}

\begin{document}

% ------------------------------------------------------------
\begin{center}
  {\large\bfseries Sensores y librerías sobre STM32}\\[2pt]
  {\small Laboratorio 3 \textbar{} Sistemas Electrónicos de Instrumentación}\\[2pt]
  {\footnotesize Pedro Villar \quad\textbar\quad 13 de noviembre de 2026}
\end{center}
\vspace{-4pt}
\hrule
\vspace{4pt}

Resolví los cuatro desafíos sobre una \textbf{STM32F4Discovery}, con un HC-SR04 y un LM35
cableados en protoboard y el acelerómetro LIS3DSH que ya viene en la placa. Acá va el lazo
principal de cada desafío con sus callbacks; el proyecto completo (CubeMX + CMake, un archivo
por desafío y la librería \texttt{lm35} propia) está en el
\href{https://github.com/PedroMVillar/Sist.Instrumentacion}{repositorio}.

% ------------------------------------------------------------
\section*{Desafío 1: HC-SR04, librerías de terceros y actuadores}

\begin{lstlisting}
while (1) {
  HAL_TIM_IC_Stop_IT(sensor_distancia.htim, sensor_distancia.IC_Channel);
  hcsr04_start_measurement(&sensor_distancia);  // dispara y vuelve enseguida
  HAL_Delay(100);                               // el eco tarda ~38 ms como mucho

  if (hcsr04_is_data_ready(&sensor_distancia)) {
    float distancia = hcsr04_get_distance_cm(&sensor_distancia);
    // apagar los tres LEDs y encender solo el de la zona actual
    if (distancia > 50.0f)
      HAL_GPIO_WritePin(LED_VERDE_GPIO_Port, LED_VERDE_Pin, GPIO_PIN_SET);
    else if (distancia >= 20.0f)
      HAL_GPIO_WritePin(LED_NARANJA_GPIO_Port, LED_NARANJA_Pin, GPIO_PIN_SET);
    else
      HAL_GPIO_WritePin(LED_ROJO_GPIO_Port, LED_ROJO_Pin, GPIO_PIN_SET);
  }
}

void HAL_TIM_IC_CaptureCallback(TIM_HandleTypeDef *htim) {
  hcsr04_tim_ic_callback(&sensor_distancia, htim);  // los flancos los procesa la librería
}
\end{lstlisting}

La librería de la cátedra hace el trabajo fino: genera el pulso de trigger y mide el ancho del
eco con el timer en Input Capture. Desde el main solo se dispara la medición, se espera y se
pregunta si hay dato. El \texttt{Stop\_IT} del principio evita que un eco perdido deje trabado
el arranque siguiente. El callback no interpreta nada: le pasa los flancos a la librería, que es
la que sabe qué hacer con ellos.

% ------------------------------------------------------------
\section*{Desafío 2: LM35, modularización y multi-instancia}

\begin{lstlisting}
LM35_Init(&lm35_1, &hadc1, ADC_CHANNEL_11);   // dos instancias lógicas
LM35_Init(&lm35_2, &hadc1, ADC_CHANNEL_11);   // sobre el mismo canal físico

while (1) {
  // Fase A: lectura bloqueante de la instancia 1
  temp_bloqueante = LM35_Read_Blocking(&lm35_1);

  // Fase B: lectura no bloqueante de la instancia 2
  LM35_Start_Conversion(&lm35_2);
  float temperatura;
  while (!LM35_Get_Temperature(&lm35_2, &temperatura));  // false = sigue convirtiendo
  temp_no_bloqueante = temperatura;
}
\end{lstlisting}

La gracia está en la librería propia: cada instancia es solo una estructura con el ADC y el canal
que usa, así que varias conviven sin pisarse y sin tocar el hardware a bajo nivel (de eso se
encarga CubeMX). La fase A espera la conversión adentro de la función; la fase B la dispara y
deja preguntar después, que es justo lo que necesita el desafío 4. Las dos temperaturas se
verifican en Live Expressions con el debugger.

% ------------------------------------------------------------
\section*{Desafío 3: el acelerómetro y el problema del bloqueo}

\begin{lstlisting}
while (1) {
  int8_t x = (int8_t)accel_leer_registro(LIS3DSH_REG_OUT_X_H);
  int8_t y = (int8_t)accel_leer_registro(LIS3DSH_REG_OUT_Y_H);
  ActualizarLedsInclinacion(x, y);             // naranja, verde y rojo

#if HEARTBEAT_BLOQUEANTE
  HAL_GPIO_TogglePin(LED_AZUL_GPIO_Port, LED_AZUL_Pin);  // el problema
  HAL_Delay(500);
#else
  if (HAL_GetTick() - millisAnterior >= 500) {           // la solución
    millisAnterior = HAL_GetTick();
    HAL_GPIO_TogglePin(LED_AZUL_GPIO_Port, LED_AZUL_Pin);
  }
#endif
}
\end{lstlisting}

Con \texttt{HAL\_Delay} el azul parpadea bien, pero la inclinación se revisa solo dos veces por
segundo y los LEDs responden a los tirones. Con \texttt{HAL\_GetTick} el loop gira libre: el azul
parpadea igual y la inclinación responde al instante. El azul queda reservado para el latido, así
que la inclinación usa los otros tres LEDs. Dato de la práctica: los ejes del chip vienen rotados
90 grados respecto de los LEDs de la placa, así que el mapeo de signos salió de probar.

% ------------------------------------------------------------
\section*{Desafío 4: el super-loop y la integración total}

\begin{lstlisting}
while (1) {
  uint32_t ahora = HAL_GetTick();

  // Tarea 1: acelerómetro, en cada vuelta
  int8_t x = (int8_t)accel_leer_registro(LIS3DSH_REG_OUT_X_H);
  int8_t y = (int8_t)accel_leer_registro(LIS3DSH_REG_OUT_Y_H);
  ActualizarLedsInclinacion(x, y, alarma);

  // Tarea 2: temperatura cada 1000 ms, fase B de la librería lm35
  if (ahora - ultimo_temperatura >= 1000) {
    ultimo_temperatura = ahora;
    LM35_Start_Conversion(&sensor_temperatura);
    conversion_en_curso = 1;
  }
  if (conversion_en_curso) {
    float temp_nueva;
    if (LM35_Get_Temperature(&sensor_temperatura, &temp_nueva)) {
      temperatura = temp_nueva;                // valor consolidado del segundo
      conversion_en_curso = 0;
    }
  }

  // Tarea 3: distancia cada 100 ms
  if (ahora - ultimo_distancia >= 100) {
    ultimo_distancia = ahora;
    if (hcsr04_is_data_ready(&sensor_distancia))
      distancia = hcsr04_get_distance_cm(&sensor_distancia);
    HAL_TIM_IC_Stop_IT(sensor_distancia.htim, sensor_distancia.IC_Channel);
    hcsr04_start_measurement(&sensor_distancia);
  }

  // Tarea 4: alarma por obstáculo a menos de 15 cm
  alarma = (distancia < 15.0f);
  if (alarma && ahora - ultimo_parpadeo >= 100) {   // 5 Hz
    ultimo_parpadeo = ahora;
    HAL_GPIO_TogglePin(LED_ROJO_GPIO_Port, LED_ROJO_Pin);
  }
}
\end{lstlisting}

No hay un solo \texttt{HAL\_Delay}: cada tarea guarda la última vez que corrió y compara contra
el tick actual, así ninguna frena a las demás. El acelerómetro se lee siempre, la temperatura usa
la fase B de \texttt{lm35} y la distancia se recoge y se vuelve a disparar cada 100 ms. La alarma
es solo una bandera: cuando hay obstáculo cerca, el rojo deja de obedecer al acelerómetro y
parpadea a 5\,Hz; al despejarse, vuelve solo a indicar inclinación.

% ------------------------------------------------------------
\section*{Conclusiones}

El laboratorio junta las tres formas de entrar datos al micro: un sensor digital que se mide con
un timer (HC-SR04), uno analógico que pasa por el ADC (LM35) y uno que habla por un bus serie
(el acelerómetro por SPI). Del lado del software, la lección fue separar capas: la librería
resuelve el sensor y el main decide qué hacer con el dato. Escribir \texttt{lm35} con instancias
y fases me ayudó a entender la librería de la cátedra, que sigue la misma idea. Y el desafío 4
cierra lo que venía del laboratorio 2: en el super-loop final cada tarea decide si le toca correr
mirando el reloj, y por eso tres sensores y una alarma conviven en un mismo hilo sin molestarse.

\end{document}
